% ============================================================================
% RELATED: HARDWARE-ACCELERATED BOOLEAN COMPUTATION
% ============================================================================
\subsection{Hardware-Accelerated Boolean Computation}
\label{sec:related_hardware_accel}

The Value architecture is designed to align with the native instruction
sets of modern parallel hardware.  Every operation in the system reduces
to a combination of bitwise logic (AND, OR, XOR, NOT), population count
(popcount), and integer comparison---instructions that execute in a
single clock cycle on contemporary CPU, GPU, and accelerator
architectures.

\paragraph{SIMD and GPU Mapping.}
Each Value is 1024 bytes, or sixteen 512-bit SIMD registers on x86
(AVX-512) and thirty-two 256-bit registers on ARM (NEON).  A single
boolean truth-table operation across an entire Value therefore requires
at most 16--32 SIMD instructions, independent of the semantic content of
the frame.  On GPU hardware, the same operation maps to a threadgroup
dispatch where each thread processes one 64-bit word of the frame.

The system provides three backend implementations:

\begin{itemize}[nosep]
  \item \textbf{CPU:} A portable implementation that iterates over the
        128 words of each frame, applying the selected truth table via
        bitwise logic and writing the result in-place.  Population count
        is computed via the hardware \texttt{popcount} intrinsic where
        available.
  \item \textbf{Apple Metal:} Compute shaders dispatch value operations
        across threadgroups, with each thread processing one 64-bit
        word.  The Metal backend is compiled from shader source at build
        time using \texttt{xcrun}.
  \item \textbf{NVIDIA CUDA:} Equivalent kernels dispatch across warps,
        with coalesced memory access patterns to maximize memory
        bandwidth utilization.
\end{itemize}

\paragraph{Deterministic Latency.}
Because every instruction is a fixed-length bitwise operation over a
fixed-size frame, execution time is data-independent: there are no
variable-length encodings, no data-dependent branches, and no
floating-point precision concerns.  This property makes worst-case
latency trivially analyzable and eliminates the tail-latency variance
that characterizes matrix-multiplication-based inference on the same
hardware.

\paragraph{Write-Optimized Access Patterns.}
The Region's multi-tier buffering structure (primary channel, spill
channel, overflow queue) is designed for write-heavy workloads.  Writes
never block: a frame is always accepted into the next available tier.
Reads are non-blocking and return \texttt{io.EOF} when the Region is
empty.  This lock-free design avoids the contention that arises in
shared-memory architectures when multiple producers and consumers
compete for the same data structure, and it ensures that the GPU
backends are never stalled waiting for a synchronization barrier.

\paragraph{Relationship to In-Memory Computing.}
The all-boolean, fixed-size nature of Value operations is particularly
well-suited to emerging in-memory computing architectures, where bitwise
operations can be performed directly within memory arrays without
transferring data to a separate processing unit.  While the current
implementation uses conventional CPU and GPU hardware, the instruction
set imposes no requirements beyond bitwise logic and population count,
making it portable to any substrate that supports these primitives.
